\partlabel{PART I --- MANUFACTURING PROBLEM AND SYSTEM MODEL}
\chapter{The Manufacturing Problem}
A polymer-compounding line converts a variable feed into a controlled melt. Its operator must distinguish a normal change of recipe or throughput from developing mechanical or thermal trouble. An increase in screw torque together with gearbox vibration raises a practical question: what changed, which evidence supports a possible cause, and which inspection should a qualified technician prioritize?

The project's example is a \textbf{Coperion ZSK co-rotating twin-screw extruder}. The running RAG catalog records ``ZSK Twin-Screw Extrusion Line 01'', model \textbf{ZSK 58}, manufacturer Coperion GmbH, application Polymer Compounding, and synthetic data provenance. This is a simulated configuration, not evidence that a physical ZSK 58 is connected. The measurement CSV identifies \texttt{ZSK-SIM-01}; the RAG metadata identifies \texttt{ZSK-001}. No verified asset mapping joins them automatically.
\section{Components and monitored behavior}
The main drive motor supplies mechanical power through the gearbox to two intermeshing screws rotating in the same direction. Screw elements convey, mix and shear material within the barrel. Heating and cooling zones regulate processing temperature. The feeder establishes input mass flow; melt-processing components and the outlet establish pressure and downstream resistance. The software stores one aggregate barrel-temperature signal, not a separately instrumented array of zones.
\fig{manufacturing}{Original schematic of the simulated line. Boxes are functional components, not a manufacturer drawing or an installed instrumentation layout.}
\section{The decision is an evidence problem}
A vibration rise can be consistent with bearing deterioration, alignment changes or lubrication problems, but RMS vibration alone cannot establish a bearing defect. High torque can reflect process resistance rather than gearbox damage. Recipe, loading, thermal conditions and maintenance history therefore matter. Coperion describes gearbox condition monitoring based on vibrations and oil condition; this motivates the example but does not validate the project's classifier~\cite{coperionmonitor}.
\verdict{The extruder is a useful integration case because interacting process and mechanical signals require contextual interpretation. The software remains decision support: it neither authorizes physical maintenance nor issues equipment-control commands.}
\source{\path{zsk-service/data/zsk_labeled_1000.csv}; \path{zsk-service/zsk/schemas.py}: \texttt{Measurement}; live RAG catalog checked 9 October 2026; \path{rag-server/rag/models.py}: \texttt{Device}.}
